\documentclass[11pt]{article}
\usepackage[margin=1in]{geometry}
\usepackage{amsmath,amssymb,amsthm}
\usepackage{xurl}
\usepackage{hyperref}
\sloppy
\let\oldus\_ \renewcommand{\_}{\oldus\allowbreak}
\newtheorem{theorem}{Theorem}
\newtheorem{lemma}[theorem]{Lemma}
\theoremstyle{definition}
\newtheorem{definition}[theorem]{Definition}
\newcommand{\qb}[2]{\genfrac{[}{]}{0pt}{}{#1}{#2}_{q}}
\newcommand{\F}{\mathbb{F}}

\title{Proof of conjecture 00000002478:\\ the q-binomial (Gauss) expansion and its finite-field counting content}
\author{Jackmeson1}
\date{2026-10-05}

\begin{document}
\maketitle

\section{The conjecture and its reading}

\textbf{Statement (English, verbatim).} Definition: The q-binomial theorem: the Gauss
expansion. Conjecture: The q-binomial theorem generalizes to the Gauss expansion, with the
counting content given by the lattices of finite fields -- the q-binomial coefficient counts
subspaces, and the expansion closes under the Gaussian lattice. (q-binomial theorem Gauss
expansion finite-field lattice)

\textbf{Chinese text.} It says, in substance, that the q-binomial theorem has a Gauss
generalization, and that its combinatorics is given by the lattices of finite fields.

\textbf{These are classical theorems, not new results.} We prove them and check the proofs in
Lean 4 with Mathlib. We read the statement as four claims. Each one is a theorem in the Lean
file.
\begin{itemize}
\item[(E)] \emph{The expansion.} As polynomials in two variables $q,t$ with integer (indeed
  natural) coefficients,
  \[\prod_{i=0}^{n-1}(1+q^it)=\sum_{k=0}^{n}q^{k(k-1)/2}\qb{n}{k}t^k .\]
  The same identity holds after evaluating $q$ and $t$ in any commutative semiring. Wikipedia
  [W] states this exact formula and names it the ``Cauchy binomial theorem''. The conjecture
  calls it the Gauss expansion.
\item[(C1)] \emph{The coefficient counts subspaces.} Let $\F$ be a finite field with $q$
  elements. The number $N_q(n,k)$ of $k$-dimensional subspaces of $\F^n$ is $\qb{n}{k}$
  evaluated at $q$.
\item[(C2)] \emph{Counting content of the expansion.} For $q=|\F|$,
  $\prod_{i<n}(1+q^it)=\sum_{k\le n}q^{k(k-1)/2}N_q(n,k)\,t^k$ in $\mathbb{N}[t]$.
\item[(P)] \emph{The expansion closes.} The coefficients and the counts satisfy the q-Pascal
  recursion $\qb{n+1}{k+1}=\qb{n}{k}+q^{k+1}\qb{n}{k+1}$. Multiplying by one more factor turns
  the order-$n$ expansion into the order-$(n+1)$ expansion exactly by this recursion; see the
  proof of (E).
\end{itemize}
\textbf{Not covered here.} We read ``closes under the Gaussian lattice'' as (P). A different
reading is \emph{interval closure}: every interval of the subspace lattice is again counted by
Gaussian binomials. That reading belongs to the subspace-count conjecture 00000002490, and it is
proved in our separate package for that conjecture. This package does not repeat it.

\textbf{Conventions.} $\F$ is any finite field with $q=|\F|\ge 2$; its prime-power form is not
used. A subspace is a Mathlib \texttt{Submodule}. Dimension is \texttt{finrank}. A count is
\texttt{Nat.card}. $k(k-1)/2$ is natural-number division, and it equals $\binom{k}{2}$
(\texttt{Nat.choose\_two\_right}). The two-variable polynomial ring is $\mathbb{N}[q][t]$, with
$q$ the inner variable and $t$ the outer one.

\section{Definitions as formalized}

\begin{definition}[\texttt{gaussBinom}]
The Gaussian binomial $\qb{n}{k}\in\mathbb{N}[q]$ is defined by the rules $\qb{n}{0}=1$,
$\qb{0}{k+1}=0$, and $\qb{n+1}{k+1}=\qb{n}{k}+q^{k+1}\qb{n}{k+1}$.
\end{definition}

This is the first ``analog of Pascal's identity'' in [W]. Wikipedia instead \emph{defines} the
coefficient by the quotient $(1-q^m)\cdots(1-q^{m-r+1})/((1-q)\cdots(1-q^r))$. We prove
(\texttt{gaussBinom\_mul\_qDen}) that
$\qb nk\prod_{i<k}(1-q^{i+1})=\prod_{i<k}(1-q^{n-i})$ holds in every commutative ring, so the
two definitions agree. For $k>n$ both sides are $0$, which matches [W] (``If r > m, this
evaluates to 0''). We also show that $\qb nk=0$ for $n<k$ (\texttt{gaussBinom\_eq\_zero\_of\_lt}).

\section{Proofs}

\begin{theorem}[(E), \texttt{qBinomial\_eval}, \texttt{qBinomial\_expansion}]
In every commutative semiring, for all $x,t$,
$\prod_{i<n}(1+x^it)=\sum_{k\le n}x^{\binom k2}\qb nk(x)\,t^k$.
\end{theorem}
\begin{proof}
The proof is by induction on $n$, for all $t$ at once; $n=0$ is trivial. Peel off the factor
$i=0$:
\[\prod_{i<n+1}(1+x^it)=(1+t)\prod_{i<n}\bigl(1+x^i(xt)\bigr)=(1+t)\sum_{k\le n}a_k\,x^kt^k,\]
where $a_k=x^{\binom k2}\qb nk(x)$ and the induction hypothesis is used at $xt$. Since
$\qb{n}{n+1}=0$, we have $\sum_{k\le n}a_kx^kt^k=1+\sum_{k\le n}a_{k+1}x^{k+1}t^{k+1}$. So the
coefficient of $t^{k+1}$ is
$a_{k+1}x^{k+1}+a_kx^k=x^{\binom{k+1}2}\bigl(\qb nk+x^{k+1}\qb n{k+1}\bigr)=x^{\binom{k+1}2}\qb{n+1}{k+1}$.
Here we used $\binom{k+1}2=\binom k2+k$ and the q-Pascal rule (P). The constant term is $1$ on
both sides. The polynomial identity follows by evaluating at $x=q$ and $t$ in
$\mathbb{N}[q][t]$.
\end{proof}

\begin{theorem}[(C1), \texttt{card\_subspaces\_eq\_gaussBinom}, \texttt{subCount\_eq}]
Let $V$ be a finite $\F$-vector space of dimension $n$. Then
$\#\{W\le V:\dim W=k\}=\qb nk(q)$.
\end{theorem}
\begin{proof}
Double count the ordered linearly independent $k$-tuples. Each one spans a $k$-dimensional
subspace $W$. The tuples that span a fixed $W$ are exactly the linearly independent $k$-tuples
of $W$ (\texttt{fibreEquiv}). Mathlib's \texttt{card\_linearIndependent} counts the linearly
independent $k$-tuples in an $m$-dimensional space as $\prod_{i<k}(q^m-q^i)$ when $k\le m$. This
gives $\prod_{i<k}(q^n-q^i)=N\cdot\prod_{i<k}(q^k-q^i)$ for $k\le n$ (\texttt{card\_li\_eq}).
By induction on $n$ via q-Pascal,
$\qb nk(x)\prod_{i<k}(x^k-x^i)=\prod_{i<k}(x^n-x^i)$ in every commutative ring
(\texttt{gaussBinom\_mul\_liProd}). Since $q\ge2$, we may cancel
$\prod_{i<k}(q^k-q^i)\neq0$ in $\mathbb{Z}$. For $k>n$, both $N$ and $\qb nk(q)$ are $0$.
\end{proof}

(C2) is (E) evaluated at $x=q$, with each $\qb nk(q)$ replaced by $N_q(n,k)$ using (C1)
(\texttt{qBinomial\_expansion\_count}). (P) for the counts is the defining recursion, evaluated
at $q$ and transported by (C1) (\texttt{subCount\_pascal}).

\section{Lean formalization}

The file is \texttt{lean/Conjecture2478/Basic.lean}, namespace \texttt{C2478}, 304 lines. Its
only import is \texttt{Mathlib}. The main theorem \texttt{conjecture\_2478} is a conjunction.
It states (E) as an identity in $\mathbb{N}[q][t]$ (Lean \texttt{Polynomial (Polynomial Nat)})
for all $n$. For every finite field $\F$, it also states (C1) for \texttt{Fin n -> F}, (C2), and
(P) for the counts. Supporting theorems: \texttt{qBinomial\_eval},
\texttt{qBinomial\_expansion}, \texttt{gaussBinom\_mul\_qDen},
\texttt{card\_subspaces\_eq\_gaussBinom}, \texttt{qBinomial\_expansion\_count} and
\texttt{subCount\_pascal}. The subspace-count part (\texttt{spanMap}, \texttt{fibreEquiv},
\texttt{card\_li\_eq}, \texttt{gaussBinom\_mul\_liProd}) is shared with our package for
00000002490 (same author).

\textbf{Axiom audit.} \texttt{lake build} succeeds. \texttt{\#print axioms} reports
\texttt{[propext, Classical.choice, Quot.sound]} for \texttt{conjecture\_2478},
\texttt{qBinomial\_expansion}, \texttt{qBinomial\_expansion\_count},
\texttt{gaussBinom\_mul\_qDen} and \texttt{card\_subspaces\_eq\_gaussBinom}. The file contains
no \texttt{sorry} and no \texttt{native\_decide}, and it adds no axioms.

\textbf{Related package.} Conjecture 00000002490 is about the subspace count itself. Its
package proves (C1), both q-Pascal rules for the counts, and the interval closure of the
subspace lattice. This package is about the expansion (E) and its counting interpretation (C2).

\section*{Sources}
[W] Wikipedia, ``Gaussian binomial coefficient'',
\url{https://en.wikipedia.org/wiki/Gaussian_binomial_coefficient} (retrieved 2026-10-05).
Section ``q-binomial theorem'': ``There is an analog of the binomial theorem for q-binomial
coefficients, known as the Cauchy binomial theorem'', followed by the displayed formula (E). The
page also gives the product-formula definition and the Pascal analogs, and states that
$\binom nk_q$ ``counts the number of k-dimensional vector subspaces of an n-dimensional vector
space over'' $\F_q$.

Mathlib: \texttt{card\_linearIndependent}, in
\path{Mathlib/LinearAlgebra/Matrix/GeneralLinearGroup/Card.lean}.

\end{document}
